% !TEX root = updates.tex
\begin{addblock}
\subsection{Readiness: maintenance, engineering casualties and oilers}\label{sec:readiness}
A year-long game needs the things that keep a real fleet from being at sea all the time.

\textbf{Maintenance:} a ship operates $M_{\text{on}}=90$ periods (six months) and then spends $M_{\text{in}}=45$
periods (three months) in maintenance at a home port, so in steady state a third of each fleet is unavailable.
While a ship is available its maintenance clock runs down, $r^{t+1}_j=r^t_j-1$. When $r^t_j\le0$ the ship is on
forced return to its nearest home port (not to the oiler). On arrival it enters maintenance, $z_j=M_{\text{in}}$, and
$z_j$ falls by one each period. When $z_j$ reaches $0$ the ship is available again with $r_j=M_{\text{on}}$ and full
endurance $d_j=D_k$. A ship with $z_j>0$ stays in port and is not part of either player's decision.

\textbf{Staggered start:} at $t=1$ a share $m_0=1/3$ of each player's ships, chosen at random, is in maintenance with
$z^1_j$ drawn uniformly from $\{1,\dots,M_{\text{in}}\}$. Every other ship has $r^1_j$ drawn uniformly from
$\{1,\dots,M_{\text{on}}\}$, so ships come due at different times and not all at once. These draws use their own
random stream, so a given seed has the same maintenance schedule for every pair of policies.

\textbf{Engineering casualties:} each day at sea a ship suffers an engineering casualty with probability $1\%$, so
in a two-day period the chance is $p_c=1-0.99^2\approx0.02$. Then $e_j=1$ and the ship is on forced return to its
nearest home port. On arrival it is repaired for $Z_c=5$ periods (ten days), $z_j=Z_c$, after which $e_j=0$ and
$d_j=D_k$. A ship that has a casualty while heading in for maintenance is repaired during the maintenance period.

\textbf{Oilers:} each player has one oiler. It is not commanded ship by ship. It follows a fixed rule that keeps it
where its side is doing well. The side's \emph{best region} is the one with the highest value-weighted control from
that side's point of view,
\[
r^{*,t}_i=\arg\max_{r\in\{\text{N},\text{C},\text{S}\}}\ \sigma_i\,\frac{\sum_{v\in r}V^i_v\,c^t_v}{\sum_{v\in r}V^i_v}.
\]
The oiler picks a zone of $r^{*,t}_i$ uniformly at random as its target, sails one edge per period toward it, and
picks a new random target when it arrives or when the best region changes. So it wanders among the zones of the
best region and follows when that region changes. Its own fuel $d^{o}_i$ falls by one per period away from port.
When $d^{o,t}_i\le\mathrm{dist}_i(o^t_i,H_i)$ it heads for its nearest home port, and one period there refills it to
$D_o=15$ (30 days). A ship that spends a full period in the oiler's zone without moving is refueled to $D_k$, as
in port. The oiler obeys the same storm and port rules as warships. It adds no strength and cannot be attacked yet.
\end{addblock}
